\documentclass[conference]{IEEEtran}

\IEEEoverridecommandlockouts
% The preceding line is only needed to identify funding in the first footnote. If that is unneeded, please comment it out.
\usepackage{cite}
\usepackage{amsmath,amssymb,amsfonts}
\usepackage{algorithmic}
\usepackage{graphicx}
\usepackage{textcomp}
\usepackage{xcolor}
\usepackage[linesnumbered,ruled]{algorithm2e}
%\usepackage{algorithm}
\usepackage{subfigure}
%\usepackage{subfig}
%\usepackage{subcaption}
\usepackage{comment}
\usepackage{framed}
\usepackage{enumitem}
\usepackage{ragged2e}
\usepackage{tikz}
\usepackage{booktabs,tabularx}
% line space issue
%\usepackage[singlespacing]{setspace} 
%\setstretch{1.01352} 
\usepackage[singlespacing]{setspace} \setstretch{1.0041841} % 12 pt/11.95pt

\pdfminorversion=4 

%\def\BibTeX{{\rm B\kern-.05em{\sc i\kern-.025em b}\kern-.08em
%    T\kern-.1667em\lower.7ex\hbox{E}\kern-.125emX}}
    
    
\graphicspath{{fig/}{images/}}
	\DeclareGraphicsExtensions{.eps,.pdf,.jpeg,.png}
	
% \newcommand{\our}{\textsc{DIRS}}
\newcommand{\our}{\textsc{DCTCP-FQ}}

\begin{document}

\title{ NAK-Broadcast Assisted Collision Resolution in IEEE 802.11 DCF Using Variable Packet Lengths
}

\author{\IEEEauthorblockN{
Haoyu Wang\IEEEauthorrefmark{1},
Xiaoqian Zhang\IEEEauthorrefmark{2}, 
Allen Yang\IEEEauthorrefmark{1},
Bo Sheng\IEEEauthorrefmark{1}
%Ningfang Mi\IEEEauthorrefmark{1}}
%\IEEEauthorblockA{\IEEEauthorrefmark{1}Department of Electrical and Computer Engineering, Northeastern University, Boston, USA
%\\Email: jia.da@husky.neu.edu, bhimani@ece.neu.edu, ningfang@ece.neu.edu 
%}
\IEEEauthorblockA{\IEEEauthorrefmark{1}Department of Computer Science, University of Massachusetts Boston, Boston, USA
%\\Email: sonnam.nguyen001@umb.edu, shengbo@cs.umb.edu 
}
\IEEEauthorblockA{\IEEEauthorrefmark{2}Department of Computer Science, University of Nebraska Omaha, Omaha, USA
}
%\IEEEauthorblockA{\IEEEauthorrefmark{3}Samsung Semiconductor Inc., San Jose, CA, USA
}
}

%  \author[1]{Danlin Jia}
%  \author[2]{Tengpeng Li}
%  \author[2]{Xiaoqian Zhang}
%  \author[1]{Li Wnag}
%  \author[2]{Bo Sheng}
%  \author[1]{Ningfang Mi}
%  \affil[1]{Department of Electrical and Computer Engineering, Northeastern University}
%  \affil[2]{Department of Computer Science, University of Massachusetts Boston}

%  \renewcommand\Authands{ and }
 
% \author{\IEEEauthorblockN{1\textsuperscript{st} Danlin Jia}
% \IEEEauthorblockA{\textit{Department of Electrical and Computer Engineering} \\
% \textit{Northeastern University}\\
% Boston, USA \\
% jia.da@husky.neu.edu}
% \and
% \IEEEauthorblockN{2\textsuperscript{nd} Given Name Surname}
% \IEEEauthorblockA{\textit{dept. name of organization (of Aff.)} \\
% \textit{name of organization (of Aff.)}\\
% City, Country \\
% email address}
% \and
% \IEEEauthorblockN{3\textsuperscript{rd} Given Name Surname}
% \IEEEauthorblockA{\textit{dept. name of organization (of Aff.)} \\
% \textit{name of organization (of Aff.)}\\
% City, Country \\
% email address}
% \and
% \IEEEauthorblockN{4\textsuperscript{th} Given Name Surname}
% \IEEEauthorblockA{\textit{dept. name of organization (of Aff.)} \\
% \textit{name of organization (of Aff.)}\\
% City, Country \\
% email address}
% \and
% \IEEEauthorblockN{5\textsuperscript{th} Given Name Surname}
% \IEEEauthorblockA{\textit{dept. name of organization (of Aff.)} \\
% \textit{name of organization (of Aff.)}\\
% City, Country \\
% email address}
% }


% \author{\IEEEauthorblockN{1\textsuperscript{st} Danlin Jia}
% \IEEEauthorblockA{\textit{Department of Electrical and Computer Engineering} \\
% \textit{Northeastern University}\\
% Boston, USA \\
% jia.da@husky.neu.edu }
% \and

% \IEEEauthorblockN{2\textsuperscript{nd} Janki Bhimani}
% \IEEEauthorblockA{\textit{Department of Electrical and Computer Engineering} \\
% \textit{Northeastern University}\\
% Boston, USA \\
% bhimani@ece.neu.edu}
% \and

% \IEEEauthorblockN{3\textsuperscript{rd} Son Nam Nguyen}
% \IEEEauthorblockA{\textit{Department of Computer Science} \\
% \textit{University of Massachusetts Boston}\\
% Boston, USA \\
% SonNam.Nguyen001@umb.edu}
% \and

% \IEEEauthorblockN{4\textsuperscript{th} Bo Sheng}
% \IEEEauthorblockA{\textit{Department of Computer Science} \\
% \textit{University of Massachusetts Boston}\\
% Boston, USA \\
% shengbo@cs.umb.edu}
% \and
% \IEEEauthorblockN{5\textsuperscript{th} Ningfang Mi}
% \IEEEauthorblockA{\textit{Department of Electrical and Computer Engineering} \\
% \textit{Northeastern University}\\
% Boston, USA \\
% ningfang@ece.neu.edu}

% }

% make the title area
\maketitle

% abstract.tex:

\begin{abstract}


  

\end{abstract}

\begin{IEEEkeywords}
placeholder
\end{IEEEkeywords}

%\vspace{-1mm}
\input{01_intro.tex}

%\vspace{-1mm}
\input{02_motivation.tex}

%\vspace{-1mm}
\section{Proposed Design}
\label{sec:proposal}



%\vspace{-1mm}
\input{04_evaluation.tex}

%\vspace{-1mm}
\input{05_relatedWork.tex}

%\vspace{-1mm}

\input{06_conclusion.tex}




% trigger a \newpage just before the given reference
% number - used to balance the columns on the last page
% adjust value as needed - may need to be readjusted if
% the document is modified later
%\IEEEtriggeratref{8}
% The "triggered" command can be changed if desired:
%\IEEEtriggercmd{\enlargethispage{-5in}}

% references section

% can use a bibliography generated by BibTeX as a .bbl file
% BibTeX documentation can be easily obtained at:
% http://www.ctan.org/tex-archive/biblio/bibtex/contrib/doc/
% The IEEEtran BibTeX style support page is at:
% http://www.michaelshell.org/tex/ieeetran/bibtex/
%\bibliographystyle{IEEEtranS}
% argument is your BibTeX string definitions and bibliography database(s)
%\bibliography{IEEEabrv,../bib/paper}
%
% <OR> manually copy in the resultant .bbl file
% set second argument of \begin to the number of references
% (used to reserve space for the reference number labels box)
% \begin{thebibliography}{1}

% \bibitem{IEEEhowto:kopka}
% H.~Kopka and P.~W. Daly, \emph{A Guide to \LaTeX}, 3rd~ed.\hskip 1em plus
%   0.5em minus 0.4em\relax Harlow, England: Addison-Wesley, 1999.

% \end{thebibliography}

\RaggedRight
\bibliographystyle{IEEEtran}
% include bibliography definition
\bibliography{IEEEcustom,bo, liwang,xiaoqian}

% \begin{thebibliography}{00}
%\bibitem{Delay scheduling} M. Zaharia et al. Delay scheduling: A simple technique for achieving locality and fairness in cluster scheduling. In EuroSys 2010
%\bibitem{Fujitsu VDI Download link} “Fujitsu VDI,” [Online]. Available: http://iotta.snia.org/tracetypes/3
%\bibitem{Fujitsu K5cloud Download link} "Fujitsu K5cloud," [Online]. Available: https://www.fujitsu.com/kz/Images
%\bibitem{Tencent CBS Download link} “Tencent Cloud Block Storage (TCB),” [Online]. Available: https://intl.cloud.tencent.com/product/cbs
%\bibitem{MATLAB} “MATLAB,” [Online]. Available: https://www.mathworks.com
%\bibitem{Williams and Mary} “Q-MAP,” [Online]. Available: https://www.wm.edu/as/computerscience/
%\bibitem{BMAP-Trace} “BMAP” [Online]. Available: https://
%\bibitem{acs Verification} “./acs verification ,” [Online]. Available: https://


%\bibitem{PCN} Cheng W, Qian K, Jiang W, Zhang T, Ren F. Re-architecting congestion management in lossless Ethernet.NSDI 2020
%\bibitem{TCP} Padhye J, Firoiu V, Towsley D, Kurose J. Modeling TCP throughput: A simple model and its empirical validation.SIGCOMM 1998.
%\bibitem{PFC} 802.1Qbb – Priority-based Flow Control.https://1.ieee802.org/dcb/802-1qbb/
%\bibitem{QCN} 802.1Qau – Congestion Notification. https://1.ieee802.org/dcb/802-1qau/
%\bibitem{ECN} Floyd S. TCP and explicit congestion notification.SIGCOMM 1994
%\bibitem{DCTCP} Alizadeh, Mohammad, et al. Data center tcp (dctcp). SIGCOMM 2010
%\bibitem{DCQCN} Zhu, Yibo, et al. Congestion control for large-scale RDMA deployments.SIGCOMM 2015.
%\bibitem{TIMELY} Mittal, Radhika, et al. TIMELY: RTT-based congestion control for the datacenter. SIGCOMM 2015.
%\bibitem{PFABRIC} Alizadeh, Mohammad, et al. pfabric: Minimal near-optimal datacenter transport.SIGCOMM 2013.
%\bibitem{FASTPASS} Perry, Jonathan, et al. Fastpass: A centralized" zero-queue" datacenter network. SIGCOMM 2014
%\bibtem{PHOST} Gao, Peter X., et al. phost: Distributed near-optimal datacenter transport over commodity network fabric. CoNEXT 2015
%\bibtem{NDP} Handley, Mark, et al. Re-architecting datacenter networks and stacks for low latency and high performance. SIGCOMM 2017
%\bibtem{HOMA} Montazeri, Behnam, et al. Homa: A receiver-driven low-latency transport protocol using network priorities. SIGCOMM 2018
% \end{thebibliography}

% that's all folks
\end{document}


